% Rozdział 4 — Stanowisko pomiarowe i metodyka (budżet: 3,5 str.)
% Źródła: uklad-pomiarowy/uklad_pomiarowy_body.tex (schemat circuitikz,
% samodzielny dokument — wstawiany makrem \inputBezPreambuly, bez zmian w pliku).
% Podpis jak w uklad-pomiarowy/rysunek_w_pracy.tex, ze ścieżką i nazwą
% dopasowaną do pracy (słownik N6, N9). Docelowo wątek rozdziału
% przenosi ten blok do własnego pliku.
\chapter{Stanowisko pomiarowe i~metodyka}
\label{ch:stanowisko}
\TODO{opis komory próżniowej, obwodu E--B--T, spektrometru, metody analizy zdjęć}

\begin{figure}[htbp]
  \centering
  \resizebox{\textwidth}{!}{\inputBezPreambuly{../uklad-pomiarowy/uklad_pomiarowy_body}}
  \caption{Schemat układu do pomiaru emisji polowej w~prototypie analizatora EDS:
    E -- emiter, B -- bramka (elektroda ekstrakcyjna), T -- target;
    $\UT$ -- zasilacz obwodu targetu, $\UB$ -- zasilacz bramki;
    $\RT$, $\RB$ -- rezystory ograniczające, V -- woltomierze
    \opracowanieWlasne}
  \label{fig:uklad-pomiarowy}
\end{figure}
